\begin{minipage}{0.65\linewidth}
    Dans un jeu de lancement sur un rail $AB$ rectiligne horizontal, un solide
    $\mathrm{(S_{0})}$ de masse $m_{0}$ et de centre d'inertie $G$ est lancé à
    partir d'une position $A$ avec une vitesse initiale $\vec{v}_{0}$
    horizontale pour atteindre une cible $\mathrm{(C)}$ qui se trouve en $B$ à
    la distance $d=AB$
    (Fig.~\ref{fig:2bac_svt_national_2025_normale_phys_ex3_1}).
    
    Le mouvement de $\mathrm{(S_{0})}$ sur le rail $AB$ se fait avec des
    frottements modélisés par une force $\vec{f}$ constante de sens opposé au
    sens du vecteur vitesse.
\end{minipage}
\hfill
\begin{minipage}{0.34\linewidth}
    \begin{figure}[H]
        \centering
        \begin{tikzpicture}[]
            \Sol{5cm}{3mm}{sol}
            \node[solide={7mm/7mm},anchor=south west] (S) at (sol.north west) {};
            \node[point,label={[above=3mm]:$A$}] (A) at (S.center) {};
            \draw[->] (A.center) -- ++(2,0)
                node[anchor=south east,inner xsep=0pt] {$\vec{v}_{0}$};
            \draw[->,thick] (A.center) -- ++(1,0)
                node[anchor=south east,inner xsep=0pt] {$\ex$};
            \node[minimum width=1mm,minimum height=1cm,anchor=south east,
                inner xsep=0pt,fill,label=right:$B$] (C) at (sol.north east) {};
            \node[anchor=south east,inner xsep=0pt,
                inner ysep=1pt] at (C.north east) {Cible $\mathrm{(C)}$};
        \end{tikzpicture}
        \caption{}
        \label{fig:2bac_svt_national_2025_normale_phys_ex3_1}
    \end{figure}
\end{minipage}
    
Pour étudier le mouvement de $G$ sur $AB$, on choisit le repère $(A,\ex)$ lié à
la Terre supposé galiléen.
    
À $t_{0}=0$~: $x_{G}=x_{A}=0$.

% \begin{donnees}
%     \begin{itemize}
%         \item $g=\qty{10}{\m\per\square\s}$~; $d=AB=\qty{1}{\m}$~;
%               $m_{0}=\qty{0.1}{\kg}$.
%     \end{itemize}
% \end{donnees}

\textbf{\textsc{Données~:}}
$g=\qty{10}{\m\per\square\s}$~; $d=AB=\qty{1}{\m}$~; $m_{0}=\qty{0.1}{\kg}$.

\begin{enumerate}
    \item En appliquant la deuxième loi de \textsc{Newton}, montrer que
          l'équation différentielle vérifiée par l'abscisse $x_{G}$ s'écrit~:
          $\frac{\mathrm{d}^{2}x_{G}}{\mathrm{d}t^{2}}=-\frac{f}{m_{0}}$.
    \item En déduire la nature du mouvement de $G$.
    \item L'équation de la vitesse $v_{G}(t)$ de $G$ s'écrit~:
          $v_{G}(t)=-1,5\cdot t+1,5$ ($\unit{\v}$).
          \begin{enumerate}
              \item Écrire l'expression numérique de l'équation horaire du
                    mouvement de $G$.
              \item Déterminer l'instant $t_{1}$ ou $\mathrm{(S_{0})}$ s'arrête.
              \item Déterminer la valeur de $f$.
              \item Le solide $\mathrm{(S_{0})}$ atteint-il la cible
                    $\mathrm{(C)}$~? Justifier.
          \end{enumerate}
    \item Déterminer, la valeur minimale $v_{0,\min}$ de la vitesse
          initiale avec laquelle $\mathrm{(S_{0})}$ doit être lancé pour qu'il
          atteigne la cible $\mathrm{(C)}$.
\end{enumerate}

